\documentclass[
    a4paper,                % DIN A4 Papiergröße
    11pt,                   % Schriftgröße 11
    oneside,                % Einseitig
    headinclude,             % Kopfzeile wird für das Layout berücksichtigt
]{scrreprt}

% Auswahl von Zeichensatz, Sprache & Schriftart
\usepackage[utf8]{inputenc}
\usepackage[ngerman]{babel}
\usepackage[default]{opensans}
\usepackage[T1]{fontenc}

% Ermöglicht die Zählung der Abbildungen und Tabellen
\usepackage[figure, table]{totalcount}
% table setup
\usepackage{longtable}
\usepackage{array}
\usepackage{ragged2e}
\usepackage{lscape}

% listings style
\usepackage{listings}
\lstset{
	tabsize=3,
	extendedchars=false,
	inputencoding=utf8,
	frame=single,
	columns=fullflexible,
	showspaces=false,                
    showstringspaces=false,
    showtabs=false,
	language=bash,
	extendedchars=true,
    basicstyle=\ttfamily,
    keywordstyle=\lst@ifdisplaystyle\color{blue}\fi,
    commentstyle=\color{gray},
	breaklines=true,             % Zeilenumbruch	
    postbreak=\mbox{\textcolor{red}{$\hookrightarrow$}\space},  %roter Pfeil bei Zeilenumbruch
	breakatwhitespace=true, %Umbruch an Leerzeichen
	breakautoindent=true,
	framexleftmargin=15pt,
	literate={\_}{}{0\discretionary{\_}{}{\_}}%Zeilenumbrüche bei Underscrore
	         {\:}{}{0\discretionary{:}{}{:}}%Zeilenumbruch bei Doppelpunkt
	         {-}{{-\allowbreak}}{1}%Zeilenumbruch bei Bindestrich
	         {\/}{}{0\discretionary{/}{}{/}}%Zeilenumbruch bei Slash
	         {~}{{\char`\ }}1
}

% Ermöglicht das Einbindden von Bildern
\usepackage{graphicx}
\graphicspath{ {./images/} }

% Ermöglicht die Forcierung der Platzierung von Floats
\usepackage{float}

%Ermöglicht If Cases mit Strings zur Differenzierung Bamberg/Duisburg-Essen
\usepackage{xstring}

% Formatierungsoptionen
% Allgemeine Seiteneinrichtung
\usepackage[
    left=3cm,
    right=3cm,
    top=2cm,
    bottom=2cm,
    headsep = \baselineskip,
    footskip = \dimexpr2\baselineskip+3mm\relax,
    includeheadfoot
]{geometry}
\usepackage[onehalfspacing]{setspace}
\emergencystretch 3em

% Formatierung von Tabellen mit richtigem Zeilenabstand usw.
\renewcommand{\arraystretch}{1.5}
\setlength{\tabcolsep}{10pt}

% Paragraphendefinition
\setlength{\parindent}{0em}
\setlength{\parskip}{.6em}

\RedeclareSectionCommand[tocindent=0pt]{section}
\RedeclareSectionCommand[tocindent=0pt]{subsection}
%\RedeclareSectionCommand[tocnumwidth=70pt]{chapter}

\renewcommand*\chaptermarkformat{\chapappifchapterprefix{\ }% 
  \thechapter.\enskip}

% Farbdefinitionen
\usepackage{xcolor}
\definecolor{vawiblue}{HTML}{004C93}
\definecolor{vawilightblue}{HTML}{4394E0}


% Kopf- und Fußzeilen
\usepackage{scrlayer-scrpage}
\clearpairofpagestyles
\newpairofpagestyles{headeronly}{%
    \ohead*{\includegraphics[scale=0.08]{vawi-logo}}
}
\ohead*{\includegraphics[scale=0.08]{vawi-logo}}
\ofoot*{\pagemark}